problem:  
Let \((\mathfrak n, \langle\cdot,\cdot\rangle_{\mathfrak n})\) be a **nilsoliton**, i.e. a nilpotent Lie algebra with metric satisfying  
\[
\mathrm{Ric}_{\mathfrak n} = c\, I + D_E^s, \quad c<0, \quad D_E^s\in \mathrm{Der}(\mathfrak n),\ D_E^s>0.
\]
A **solvsoliton extension** of \((\mathfrak n, D_E^s)\) is a solvable Lie algebra of the form  
\[
\mathfrak s = \mathfrak a \ltimes \mathfrak n, \qquad [\mathfrak a,\mathfrak a]=0,
\]
where each \(A\in \mathfrak a\) acts on \(\mathfrak n\) by a symmetric derivation commuting with \(D_E^s\).  
Let  
\[
\mathfrak C(\mathfrak n, D_E^s)
= 
\{\text{abelian subalgebras } \mathfrak a \subset \mathrm{Der}_{\mathrm{sym}}(\mathfrak n)
\mid [A,D_E^s]=0 \text{ for all } A\in \mathfrak a\},
\]
and denote by \(\mathcal{M}_{\mathrm{sol}}(\mathfrak n, D_E^s)\) the moduli space of solvsoliton metrics arising from such extensions, up to scaling and isometry.

**Question:**  
Is the moduli space \(\mathcal{M}_{\mathrm{sol}}(\mathfrak n, D_E^s)\) determined, up to diffeomorphism, by the *conjugacy class* of the abelian subalgebra \(\mathfrak a\) inside the symmetric derivation algebra \(\mathrm{Der}_{\mathrm{sym}}(\mathfrak n)\)?  
Equivalently, is the natural map  
\[
\frac{
   \{\text{commuting symmetric derivation subalgebras } \mathfrak a \subset \mathrm{Der}(\mathfrak n)\text{ commuting with }D_E^s\}
}{
   \mathrm{Aut}(\mathfrak n)
}
\;\longrightarrow\;
\mathcal{M}_{\mathrm{sol}}(\mathfrak n, D_E^s)
\]
injective?  

If not, characterize the equivalence relation on commuting derivation subalgebras that produces distinct solvsoliton structures but geometrically isometric metrics.  
In addition, does every **continuous family** of inequivalent abelian derivation subalgebras correspond to a continuous deformation in \(\mathcal{M}_{\mathrm{sol}}(\mathfrak n, D_E^s)\), or are these moduli discretized by global geometric constraints?  

Why is it a "good" problem:  
This formulation focuses on a mathematically precise and deep issue: **the algebraic moduli of solvsolitons derived from a fixed nilsoliton base**. It directly addresses the classification problem at the heart of homogeneous Ricci soliton theory—the relationship between derivation‑algebra embeddings and geometric uniqueness. Determining whether conjugacy classes of abelian symmetric derivation subalgebras completely parameterize solvsoliton metrics would clarify to what extent solvsoliton geometry is encoded in algebraic data. A positive answer would reveal a sharp rigidity principle; a negative one would uncover a hitherto unknown geometric moduli phenomenon within homogeneous Ricci soliton spaces. The problem elegantly bridges Lie algebra theory, representation theory, and geometric analysis, embodying the “narrow entry, deep consequence” criterion for a first‑rate research problem.